% Generated draft: requires mathematical and source review.
\documentclass[11pt]{article}
\usepackage[margin=25mm]{geometry}
\usepackage{fontspec}
\setmainfont{Latin Modern Roman}
\usepackage{amsmath,amssymb,mathtools,unicode-math}
\setmathfont{Latin Modern Math}
\usepackage{enumitem,needspace,xurl,hyperref}
\setlength{\parindent}{0pt}
\setlength{\parskip}{5pt}
\setlength{\emergencystretch}{3em}
\newcommand{\N}{\mathbb{N}}
\newcommand{\Z}{\mathbb{Z}}
\newcommand{\Q}{\mathbb{Q}}
\newcommand{\R}{\mathbb{R}}
\newcommand{\C}{\mathbb{C}}
\newcommand{\F}{\mathbb{F}}
\newcommand{\A}{\mathbb{A}}
\newcommand{\PP}{\mathbb{P}}
\newcommand{\E}{\mathbb{E}}
\newcommand{\eps}{\varepsilon}
\newcommand{\dd}{\,\mathrm{d}}
\begin{document}
\section*{UnsolvedMath ID: 1101305}
\textbf{Source problem code:} catalogue entry
\subsection{Definitions and mathematical statement}
For a commutative ring $R$ with identity, let
\[
\operatorname{SL}_2(R)=\left\{\begin{pmatrix}a&b\\c&d\end{pmatrix}:a,b,c,d\in R,\ ad-bc=1\right\}.
\]
A group is finitely generated if some finite subset generates it as an abstract group, allowing products and inverses. Write
\[
R_1=\Z[X,X^{-1}],\qquad
R_{2,F}=F[X,X^{-1},Y,Y^{-1}],
\]
where $X,Y$ are commuting indeterminates and $F$ is a field. These Laurent polynomial rings consist of finite sums with integral, respectively $F$-valued, coefficients and arbitrary integer exponents.

The source asks:
\begin{enumerate}
\item Is $\operatorname{SL}_2(R_1)$ finitely generated?
\item For a field $F$, is $\operatorname{SL}_2(R_{2,F})$ finitely generated? The family question must retain its dependence on $F$, rather than silently restrict the field.
\end{enumerate}
Equivalently, each individual finite-generation question asks whether there exist $r\ge0$ and matrices $A_1,\ldots,A_r$ in the indicated group such that every group element is a finite product of the $A_i^{\pm1}$.

The supplied August 2026 research summary reports that the first question, and the second question for $F=\F_q$ with $q$ a prime power, remain open. It separately claims a negative answer for every infinite field $F$, by a finitely-generated-ring argument referred to as ``Work done'', Section 1, which is not included. Accordingly the finite-field case is the remaining open family in that supplied assessment, while the original statement itself includes all fields. No verification or new proof of the infinite-field claim is supplied here.

The summary also mentions reductions associated with Cohn, Suslin, Bachmuth--Mochizuki, and Chu, as well as non-finite-presentation and homological obstructions. These are different from non-finite-generation and do not by themselves answer the questions. Their exact groups, coefficient systems, and hypotheses are missing. Its elementary-generation notation $\operatorname{GE}_2$ is undefined in the record; the displayed implication ``$\operatorname{GE}_2\Rightarrow$ finitely generated'' is retained only as a reported proposed connection in these settings, not as a supplied equivalence or a criterion proved here.
\subsection{Short English statement}
Are the determinant-one two-by-two matrix groups over the one-variable integral and two-variable field Laurent polynomial rings finitely generated? The supplied August 2026 summary leaves the integral ring and all finite-field cases open, while claiming infinite fields give negative answers. The cited obstruction and elementary-generation results are incomplete and do not replace a finite-generation answer.
\subsection{Sources}
\begin{itemize}
\item Local source registry: \url{data/unsolvedmath/problems.json}, permanent UnsolvedMath ID $1101305$, source code catalogue entry, including the supplied PARTIAL-PROGRESS research summary. Dataset attribution: Ulam AI and UnsolvedMath Contributors; original-source attribution and rights are retained.
\item Mladen Bestvina, \emph{Questions in Geometric Group Theory}, source-list date 2004, Question 13.5, attributed in the supplied transcription to ``Seymour Bachmut''; supplied metadata identifies PDF page 24: \url{https://www.math.utah.edu/~bestvina/eprints/questions-updated.pdf}. The record reports manual transcription and July 29, 2026 access, with current-status and rights review flags. The attribution's spelling is not silently corrected.
\item Abramenko's Banff 2022 problem, Questions $1'$ and 2, and Zaremsky's research statement, as cited by the supplied summary for the reported open status. No precise bibliographic links or full statements are supplied.
\item Cohn, Suslin, Bachmuth--Mochizuki, and Chu are named in the summary's claimed reduction chain, with partial-results dates described collectively as 2004--2015. No complete citations, exact homological coefficients, or theorem statements are provided. The referenced ``Work done'', Section 1, is absent from the record.
\end{itemize}
\end{document}
